%% Dirichlet encoding of a constructed residue sequence.

\providecommand{\cA}{\mathcal A}
\providecommand{\Vir}{\mathrm{Vir}}
\providecommand{\Lsh}{L^{\mathrm{sh}}}
\providecommand{\ClaimStatusProvedHere}{}
\providecommand{\ClaimStatusConditional}{}
\providecommand{\ClaimStatusDefinitional}{}
\providecommand{\ClaimStatusOpen}{}

\chapter{Dirichlet Coordinates for the Shadow Sequence}
\label{ch:shadow-L-function-platonic}

\section*{Scope}

A Dirichlet series packages a sequence.  In this chapter the sequence
is the scalar output of an ordered residue package
\(\mathsf H_{\mathrm{res}}\).  Analytic continuation belongs to a
growth package \(\mathsf H_{\mathrm{asymp}}\); motivic realization
belongs to \(\mathsf H_{\mathrm{mot}}\); Bernoulli congruences belong
to an integral comparison package \(\mathsf H_{\mathrm{Kum}}\).
These packages keep the formal, analytic, motivic, and arithmetic
statements on their natural mathematical surfaces.

\section{Formal Dirichlet encoding}
\label{sec:shL-setup}

\begin{definition}[Shadow Dirichlet series]
\label{def:shadow-L-series}
\ClaimStatusDefinitional
Let \(\mathsf H_{\mathrm{res}}(\cA)\) define scalar coefficients
\(S_r(\cA;\mathsf H_{\mathrm{res}})\) for \(r\ge2\).  Their formal
Dirichlet series is
\[
 \Lsh(\cA;s)=\sum_{r\ge2}
 S_r(\cA;\mathsf H_{\mathrm{res}})r^{-s}.
\]
Formally, this is the coefficient sequence expressed in the Dirichlet
basis \(\{r^{-s}\}_{r\ge2}\).
\end{definition}

\begin{remark}[Three indices]
\label{rem:shL-two-reindexings}
The residue arity is \(r\), the Dirichlet variable is \(s\), and a
genus comparison uses \(g\).  A bridge between \(r\) and \(g\) is a
comparison morphism, while evaluation in \(s\) is an analytic
operation on the full sequence.
\end{remark}

\begin{proposition}[Initial half-plane]
\label{prop:shL-convergence-half-plane}
\ClaimStatusProvedHere
If \(|S_r|\le Cr^N\), then \(\Lsh(\cA;s)\) converges absolutely and
locally uniformly for \(\operatorname{Re}s>N+1\).  It is holomorphic
there and determines its coefficients uniquely.
\end{proposition}

\begin{proof}
Absolute convergence follows from comparison with
\(C\sum_{r\ge2}r^{N-\operatorname{Re}s}\).  Uniqueness is the standard
uniqueness theorem for Dirichlet series in a common half-plane.
\end{proof}

\begin{definition}[Mellin--Lerch package]
\label{def:hurwitz-lerch-shadow}
\ClaimStatusDefinitional
The package \(\mathsf H_{\mathrm{asymp}}\) consists of a polynomial
growth bound and a small-\(t\) polyhomogeneous expansion for
\[
 G_{\cA}(t)=\sum_{r\ge2}S_r(\cA;\mathsf H_{\mathrm{res}})e^{-rt},
\]
with remainder estimates stable under Mellin transform.  An equivalent
Hurwitz--Lerch presentation may be used when the coefficients admit a
normally convergent finite or countable decomposition into shifted
rational kernels.
\end{definition}

\begin{remark}[Classical continuation mechanism]
\label{rem:hurwitz-lerch-analytic-continuation}
The Mellin transform of each power-log term in the expansion of
\(G_{\cA}(t)\) is meromorphic.  The remainder estimates control the
continued integral on vertical strips.  A Hurwitz--Lerch decomposition
gives the same continuation componentwise.
\end{remark}

\section{Analytic continuation}
\label{sec:shL-continuation}

\begin{theorem}[Meromorphic continuation from
\(\mathsf H_{\mathrm{asymp}}\)]
\label{thm:shadow-L-analytic-continuation}
\ClaimStatusConditional
\emph{Type signature:}
\textup{(Open quadrant, residue Dirichlet presentation, Beilinson
levels \(2\to5\),
\(\mathsf H_{\mathrm{res}}+\mathsf H_{\mathrm{asymp}}\)).}
Under the Mellin--Lerch package, \(\Lsh(\cA;s)\) extends
meromorphically to \(\mathbb C\).  Its polar part is obtained from the
Mellin transforms of the power-log singular terms of
\(G_{\cA}(t)\), followed by multiplication by \(\Gamma(s)^{-1}\).
This final factor records the cancellations at the nonpositive
integers.
\end{theorem}

\begin{proof}
In the initial half-plane,
\[
 \Gamma(s)\Lsh(\cA;s)=
 \int_0^\infty G_{\cA}(t)t^{s-1}\,dt.
\]
Split the integral at \(t=1\), subtract finitely many terms of the
polyhomogeneous expansion at the origin, and integrate those terms
explicitly.  The controlled remainder is holomorphic in the enlarged
strip.  Iteration gives the continuation and its polar data.
\end{proof}

\begin{remark}[The weight-one motivic line]
\label{rem:shL-pole-at-s-1-motivic}
Under \(\mathsf H_{\mathrm{mot}}\), a residue assigned motivic weight
one lands in the weight-one piece of the chosen motivic period algebra.
For mixed Tate motives over \(\mathbb Z\), this piece is zero.  A
weight-zero boundary residue belongs to a distinct summand of the
Mellin model.
\end{remark}

\section{Special values}
\label{sec:shL-trivial-zeros}

\begin{theorem}[Negative-integer values]
\label{thm:trivial-zeros-structure}
\ClaimStatusConditional
Under \(\mathsf H_{\mathrm{asymp}}\), every value
\(\Lsh(\cA;-n)\), \(n\ge1\), is the evaluation of the regularized
Mellin integral.  The small-\(t\) expansion determines its
subtractions, and the holomorphic remainder records the full kernel
\(G_{\cA}\).  For the unshifted constant component,
\[
 \sum_{r\ge2}^{\mathrm{reg}}r^n
 =\zeta(-n)-1=-\frac{B_{n+1}}{n+1}-1.
\]
Shifted Lerch components contribute their corresponding Bernoulli
polynomials.
\end{theorem}

\begin{proof}
Evaluate the meromorphic Mellin continuation of
Theorem~\ref{thm:shadow-L-analytic-continuation} at \(s=-n\).  The
displayed component is the classical special-value formula for
\(\zeta(-n)\), with its \(r=1\) term removed.
\end{proof}

\begin{remark}[Functional-equation data]
\label{rem:contrast-with-riemann}
Riemann's trivial zeros follow from the completed zeta functional
equation.  A comparable zero theorem for \(\Lsh\) is governed by a
completed functional equation supplied as additional analytic data.
\end{remark}

\section{The genus comparison}
\label{sec:shL-correct-bridge}

\begin{theorem}[Genus-slot comparison]
\label{thm:Fg-from-L-sh-correctly}
\ClaimStatusConditional
Assume a chain-level shadow--Feynman comparison
\(\mathsf H_{\mathrm{gen}}\) whose arity map sends genus \(g\ge2\)
to \(r=2g-2\).  Then
\[
 \Pi_g^{\mathrm{sh}}\Lsh
 :=[r=2g-2]\Lsh=S_{2g-2}(\cA;\mathsf H_{\mathrm{res}})
\]
is the scalar residue coordinate entering the genus-\(g\) graph sum.
Theorem~D determines its uniform-weight contribution, and the
multi-weight theorem supplies the mixed-channel correction.
\end{theorem}

\begin{proof}
The first identity is coefficient extraction.  The remaining
identification is precisely the arity compatibility and graph-weight
compatibility contained in \(\mathsf H_{\mathrm{gen}}\).
\end{proof}

\begin{remark}[Bernoulli values and genus slots]
\label{rem:shL-why-naive-seemed-plausible}
Both Verlinde asymptotics and negative zeta values contain Bernoulli
numbers.  The former is a coefficient asymptotic in a moduli problem;
the latter is a regularized sum over every Dirichlet coefficient.  The
package \(\mathsf H_{\mathrm{gen}}\) states the precise bridge used in
the manuscript.
\end{remark}

\begin{remark}[Dirichlet coordinates]
\label{rem:shL-positive-content-of-correction}
The series \(\Lsh\) is a coordinate system for a constructed scalar
sequence.  Coefficient projection studies individual arities;
analytic continuation studies the global growth of the sequence.
\end{remark}

\section{Bernoulli--Kummer comparison}
\label{sec:shL-kummer}

\begin{theorem}[Kummer witnesses on the Verlinde lane]
\label{thm:kummer-congruence-prediction}
\ClaimStatusConditional
Assume \(\mathsf H_{\mathrm{Kum}}\), including the integral
normalization that identifies the relevant genus slot with the leading
coefficient of the \(\mathrm{SL}_2\) Verlinde polynomial.  Then
\[
 [n^{2(g-2)}]R_{g-2}(n^2)
 =(-1)^g\frac{2^{g-1}B_{2g-2}}{(2g-2)!},
\]
and the primes \(691\) and \(3617\) occur in the corresponding
normalized numerators at \(g=7\) and \(g=9\).
\end{theorem}

\begin{proof}
The Bernoulli coefficient is the exact cosecant-sum calculation of
Chapter~\ref{ch:z-g-kummer-bernoulli}.  The package
\(\mathsf H_{\mathrm{Kum}}\) carries it to the stated integral genus
slots; \(B_{12}=-691/2730\) and \(B_{16}=-3617/510\) give the two
witnesses.
\end{proof}

\begin{remark}[Arithmetic verification]
\label{rem:shL-falsification-test}
Expansion of the exact cosecant formula for \(Z_7\) and \(Z_9\),
followed by reduction to primitive integral numerators, verifies the
normalization component of \(\mathsf H_{\mathrm{Kum}}\).
\end{remark}

\begin{remark}[Two comparison routes]
\label{rem:shL-classical-route}
The classical route uses cosecant sums and Kummer congruences.  The
motivic route uses a motivic residue realization and its integral
period lattice.  Agreement of the two routes is a theorem inside
\(\mathsf H_{\mathrm{Kum}}\).
\end{remark}

\section{Motivic Dirichlet coordinates}
\label{sec:shL-motivic-upgrade}

\begin{corollary}[Motivic continuation]
\label{cor:shadow-L-motivic-upgrade}
\ClaimStatusConditional
Assume \(\mathsf H_{\mathrm{res}}+\mathsf H_{\mathrm{asymp}}
+\mathsf H_{\mathrm{mot}}\), with weight-compatible Mellin
subtractions.  Then poles, residues, and special values decompose by
motivic weight; the period map sends this decomposition to the
numerical continuation of \(\Lsh\).
\end{corollary}

\begin{proof}
Perform the Mellin subtraction in the graded motivic coefficient
algebra.  Each subtraction and remainder estimate respects weight by
hypothesis, and period realization commutes with the construction.
\end{proof}

\begin{remark}[Period realization]
\label{rem:shL-motivic-period-map-collapse}
The motivic decomposition distinguishes positive-weight periods from
weight-zero regularization constants.  Period realization preserves
that provenance while mapping both to numerical scalars.
\end{remark}
